(1) Michelson两臂的光都往返传播一次，因此两回波的光学路程分别为 \(2L_1\)、\(2L_2\)（忽略固定分束镜相位）。相对光程差 \(\Delta\ell=2(L_1-L_2)\)，相位差为 \(\Delta\phi=2\pi\Delta\ell/\lambda_0=4\pi(L_1-L_2)/\lambda_0\)。
(2) 若两束在探测器处光强分别为 \(I_1,I_2\)，复场相加的平均强度是 \(I=I_1+I_2+2\sqrt{I_1I_2}\cos(\Delta\phi+\phi_0)\)，其中 \(\phi_0\) 为固定光学元件相位。若两光强相等为 \(I_0\)，则 \(I=2I_0[1+\cos(\Delta\phi+\phi_0)]\)，具有清晰明暗周期。
(3) 一臂镜面平移 \(\Delta x\) 使往返光程变化 \(2\Delta x\)，故相对相位变化 \(\Delta(\Delta\phi)=2\pi(2\Delta x)/\lambda_0=4\pi\Delta x/\lambda_0\)。镜面移动半个波长会使相位变化 \(2\pi\)，探测强度经历完整一个周期。
(4) 若统计到 \(N\) 个完整条纹周期经过固定探测位置，则 \(N=|\Delta(\Delta\phi)|/(2\pi)=2|\Delta x|/\lambda_0\)，所以 \(|\Delta x|=N\lambda_0/2\)。单一强度余弦信号对相位正负改变对称，因此若要判断移动方向还需正交相位通道或其他额外测量。
